
This paper presents a computational analysis of stylistic preference drift in two RLHF annotation datasets using the Alignment Asymmetry Index (AAI), a composite instrument for measuring directional Human→AI adaptation in preference labels.

The primary finding is that verbosity preference weighting reverses sign across annotation strata in HH-RLHF (Δβ_L = +0.080, non-overlapping 95\% bootstrap CIs; 53,600 rows per stratum). This is interpreted as a population-level directional association, subject to the constraint that the available dataset lacks genuine temporal metadata or annotator identifiers.

A secondary between-group analysis on WebGPT finds statistically significant between-group differences in AI-typicality projection scores (tercile F = 33.226, p ≈ 8.3×10⁻⁹) with confirmed discriminant validity (placebo empirical p = 0.48), but does not meet the pre-registered within-annotator effect-size criterion. Two of five sub-hypotheses were not executed; one of three executed sub-hypotheses met its full pre-registered criteria; one met MUST_WORK gate criteria on code executability but not scientific effect size; and one is recorded as PARTIAL.

The findings motivate drift-aware quality auditing for multi-round RLHF annotation datasets and provide an empirically grounded specification of the data requirements for confirming individual annotator causal adaptation. The highest-priority follow-on is completing H-M3: testing whether the verbosity upweighting observable in later annotation strata propagates to reward model scoring behavior.

